% Two ways to combine 1, 2, 3 and 4 using pairs, for Section 2.2.
% Redrawn from upstream's figure 2.3.
% Shared styles for the box-and-pointer figures of Chapter 2.
%
% \input at the top of each such figure, so that the cells, leaves, pointers
% and nil strokes are drawn identically wherever they appear. A pair is ONE
% rounded rectangle with a divider through it, never two abutting boxes:
% rounding two boxes would round their inner edges too, which is not how the
% original's cells look.
\tikzset{
  pair/.style  = {draw, rounded corners=1.2pt, fill=black!5, minimum width=16mm,
                  minimum height=7mm, inner sep=0pt, outer sep=0pt},
  half/.style  = {minimum width=8mm, minimum height=7mm,
                  inner sep=0pt, outer sep=0pt},
  leaf/.style  = {draw, rounded corners=1.2pt, fill=white, minimum width=8mm,
                  minimum height=7mm, inner sep=0pt, outer sep=0pt,
                  font=\ttfamily\footnotesize},
  dot/.style   = {circle, fill, inner sep=1.1pt},
  ptr/.style   = {-{Latex[length=1.6mm, width=1.4mm]}, thin},
  plabel/.style = {font=\ttfamily\scriptsize, inner sep=1pt},
}

\begin{tikzpicture}[y = -1mm, x = 1mm]

% --- left: (cons (cons 1 2) (cons 3 4)) ---
\foreach \n/\x/\y in {R/10/0, A/6/17, B/40/0} {
  \node[half] (\n l) at (\x, \y) {};
  \node[half] (\n r) at (\x+8, \y) {};
  \node[pair] (\n)   at (\x+4, \y) {};
  \draw (\n l.north east) -- (\n l.south east);
  \node[dot] (\n d)  at (\x, \y) {};
  \node[dot] (\n rd) at (\x+8, \y) {};
}
\node[leaf] (L1) at (2, 34) {1};
\node[leaf] (L2) at (16, 34) {2};
\node[leaf] (L3) at (38, 17) {3};
\node[leaf] (L4) at (52, 17) {4};
\draw[ptr] (Rd)  -- (A.north);
\draw[ptr] (Rrd) -- (Bl.west);
\draw[ptr] (Ad)  -- (L1.north);
\draw[ptr] (Ard) -- (L2.north);
\draw[ptr] (Bd)  -- (L3.north);
\draw[ptr] (Brd) -- (L4.north);
\node[plabel] at (26, 46) {(cons (cons 1 2) (cons 3 4))};

% --- right: (cons (cons 1 (cons 2 3)) 4) ---
\foreach \n/\x/\y in {S/92/0, C/88/17, D/112/17} {
  \node[half] (\n l) at (\x, \y) {};
  \node[half] (\n r) at (\x+8, \y) {};
  \node[pair] (\n)   at (\x+4, \y) {};
  \draw (\n l.north east) -- (\n l.south east);
  \node[dot] (\n d)  at (\x, \y) {};
  \node[dot] (\n rd) at (\x+8, \y) {};
}
\node[leaf] (M4) at (120, 0) {4};
\node[leaf] (M1) at (86, 34) {1};
\node[leaf] (M2) at (110, 34) {2};
\node[leaf] (M3) at (124, 34) {3};
\draw[ptr] (Sd)  -- (C.north);
\draw[ptr] (Srd) -- (M4.west);
\draw[ptr] (Cd)  -- (M1.north);
\draw[ptr] (Crd) -- (Dl.west);
\draw[ptr] (Dd)  -- (M2.north);
\draw[ptr] (Drd) -- (M3.north);
\node[plabel] at (106, 46) {(cons (cons 1 (cons 2 3)) 4)};

\end{tikzpicture}
